\section{Object life cycle}
In \gretlpy, datasets and models (for example, \texttt{d} as a
\classpydatasetnosp instance and \texttt{m} as a
\classpymodelnosp instance) are jointly reference-managed by
\python and \lgretl. As in CPython itself, object lifetime is
governed by reference counting, with additional internal
references created during evaluation and display. Understanding
the distinction between \emph{logical deletion} and
\emph{physical destruction} is therefore essential.

\paragraph{Dataset ownership}
A \classpydataset object is logically owned by \python, but it may be
referenced indirectly by:
\begin{itemize}
	\item estimated models,
	\item cached or printed model summaries,
	\item temporary \python objects created during evaluation
		or display (e.g.\ \pythcode{print(model)} or interactive output).
\end{itemize}

Consequently:
\begin{itemize}
	\item \pythcode{del d} removes the name \texttt{d} from
		the current \python namespace,
	\item but does not necessarily destroy the underlying dataset.
\end{itemize}
The dataset is physically destroyed only when its reference count
drops to zero, i.e.\ when \emph{all} direct and indirect references
have been released.

\paragraph{Printing, evaluation, and object lifetime}
Accessing elements of \texttt{m} (such as \pythcode{\_\_repr\_\_} or
implicit interactive display) creates temporary \python objects that
may \textbf{retain references to the model summary} or other
model-related state. These objects can, in turn, hold indirect
references to the underlying dataset.

As a result, \pythcode{del d} may \textbf{not immediately trigger
	dataset destruction} if such temporary objects are still alive.
From the user's perspective, this may appear as delayed cleanup,
but it is a direct consequence of CPython's evaluation and
memory-management model.

This behavior is \textbf{intentional, well-defined, and unavoidable}
in \python, and \gretlpy follows these semantics strictly.\newline

\noindent\textbf{Example~1 (explicit printing):}
\begin{lstlisting}[style=pythonstylefoot]
import gretl

d = gretl.get_data("bjg")
m = gretl.ols("1 0 2").fit()

print(m)
del m
del d

gretl.show_data_status()
\end{lstlisting}\index{show\_data\_status()}
No dataset available.\newline

\noindent\textbf{Example~2 (implicit interactive display):}
\begin{lstlisting}[style=pythonstylefoot]
import gretl

d = gretl.get_data("bjg")
m = gretl.ols("1 0 2").fit()

m              # creates temporary references
del m
del d

gretl.show_data_status()
\end{lstlisting}\index{show\_data\_status()}
The dataset may still appear as registered for some time.
